% Discussion Section v2 — 2026-05-11
% Based on pre_draft_2026-05-10.md (ABT revised 2026-05-11)
% Outline: ¶1 gap filled · ¶2 LOCO–behavior · ¶3 literature+divergence
%          ¶4 detection/correction falsifier · ¶5 limitations · ¶6 field impact

\section{Discussion}
\label{sec:discussion}

In these two CVD cases, the selective impairment of continuous hue
interpolation (LOCO) with preserved categorical discrimination (LORO)
in hV4 identifies the cortical site at which colour distortion is
correctable. Both CVD
participants fell near floor in hV4 LOCO while maintaining
discrimination accuracy within the HC range, confirming that the
representational substrate for a correction filter is present but
geometrically distorted. A 2-component angular-dilation model captures
this distortion field with exact pre-image guarantees for all 8 hues at
both severity levels (Sub-08: mean $|\delta| = 46.3^\circ$; Sub-09:
mean $|\delta| = 20.1^\circ$), yielding individually structured
stimulus-space correction vectors from a single parameterisation family.

The dissociation between LOCO interpolation capacity and SRM $\Delta$RDM
geometric distance --- $100\%$ vs.\ $33\%$ concordance with same-day JND
thresholds --- establishes which cortical quantity is the operative
corrective target. Representational geometry (RDM) reflects the metric
structure of neural colour space; LOCO interpolation reflects whether
that structure supports functional prediction of withheld hues. That
only LOCO predicts behavioural colour confusions is consistent with
\citeA{brouwer2009}, who showed that hV4 --- but not early visual areas
--- supports interpolation to novel hues in typical observers. The
present results extend this dissociation to these two CVD cases: the same
interpolation failure that predicts JND confusion is what the correction
filter must address. Geometric distance metrics alone
are insufficient to identify it.

The S-cone axis component ($\beta_s$) of the fitted distortion field is
structurally consistent with prior evidence for cortical colour
compensation in anomalous trichromats. \citeA{emery2021} reported that
CVD perceptual errors concentrate near S-cone confusion loci in
behavioural hue-scaling tasks, and \citeA{tregillus2021} showed that
V2/V3 BOLD responses in anomalous trichromats include a cortical gain
component ($g$) that partially compensates the retinal cone shift.
The present 2-component model extends these findings to the individual
hV4 level: the residual distortion that survives V2/V3 compensation
reaches hV4 as an individually structured pattern. The correction
vectors for Sub-08 (deutan) and Sub-09 (protan) are mutually divergent
--- cosine similarity $< 0$ across hue pairs --- demonstrating directly
that a population-average retinal model cannot simultaneously serve both
CVD individuals. Where Machado-based correction predicts a shared
spectral shift for a given CVD subtype, the cortical residual is
individually signed and of opposite direction between the two cases.

All distortion models --- 1-DOF Machado, 2-DOF R+C, and 2-component ---
identify the same two participants as requiring correction, but their
prescribed correction vectors diverge in direction (cosine similarity $< 0$;
Appendix~\ref{app:retinal_family}). This detection--correction dissociation
constitutes the study's primary falsifiable prediction. In a 2AFC
experiment comparing the 2-component and retinal-family filter
conditions for each participant's most LOCO-vulnerable hue pairs, only
the cortically-grounded 2-component filter should significantly reduce
JND. If the retinal-family filter matches or exceeds the 2-component
filter in reducing discrimination thresholds, the cortical distortion
account would require revision. This preregistered comparison directly tests whether cortical or retinal
distortion is the corrective target for CVD colour correction.

Three considerations bound the proof-of-concept scope of these
findings. First, the CVD sample is $N = 2$; single-case statistics
\cite{crawford1998} permit individual-level inference but population-level
claims about deutan or protan subtypes require replication with larger
samples. Second, the confusion-axis component ($\beta_c$) is
statistically significant for Sub-08 only. For Sub-09, the bootstrap
confidence interval for $\beta_c$ includes zero, and the 2-component
model reduces to an effectively 1-component fit. The \emph{2-component}
label should not be read as confirming a two-mechanism account for both
participants. Third, case-level specificity rests on three converging
anchors --- cross-criterion significance, physiological fit magnitude,
and S-cone axis grounding --- rather than on permutation significance
alone, because the pipeline reaches nominal $p < 0.05$ in $71\%$ of HC
model--subject tests (Supplementary~\ref{app:hc_specificity}). Behavioural filter
validation remains the registered criterion for cortical-level
specificity.

This pipeline constitutes a proof-of-concept that individual cortical
colour geometry, measured via fMRI LOCO decoding and inverted through an
angular-dilation model, can ground a colour correction that is
architecturally distinct from population-average retinal approaches.
Where Machado-type correction applies a fixed spectral shift to all
individuals of a given CVD subtype, the present approach encodes each
individual's hV4 distortion field and inverts it exactly --- without
assuming population homogeneity. Phase~3 will test whether this cortical-level correction reduces
discrimination thresholds where retinal-level correction does not.
